% =========================================================
% Nederpelt & Geuvers, Type Theory and Formal Proof (2014)
% Section 1.3 -- Lambda-terms
% Notes by William Sollers, typeset in the AIM node format.
%
% Page numbers: printed page = PDF page - 29.
% Slots marked \pending{...} are yours to write.
% =========================================================
\documentclass[11pt]{article}
\usepackage{../aim-notes}

\newcommand{\Lam}{\Lambda}
\newcommand{\Var}{V}

\title{Nederpelt \& Geuvers \S1.3: Lambda-terms}
\author{William Sollers}
\date{Notes, draft 2}

\begin{document}
\maketitle

% ---------------------------------------------------------
\section{The set of $\lambda$-terms}

\begin{assumption}[Variables]\label{ass:variables}
\source{N\&G \S1.3, p.~4 (PDF p.~33).}
There is an infinite set $\Var = \{x, y, z, \dots\}$ of variables.
\end{assumption}

\begin{remark}[Open]
The book does not say whether $\Var$ is countable; the listing suggests it.
\textbf{Verify} against the page before relying on it.
\end{remark}

\begin{definition}[The set $\Lam$ of all $\lambda$-terms]\label{def:lambda-terms}
\uses{ass:variables}
\source{N\&G Definition~1.3.2, p.~4 (PDF p.~33).}
\begin{enumerate}[label=(\arabic*)]
  \item (Variable) If $u \in \Var$, then $u \in \Lam$.
  \item (Application) If $M, N \in \Lam$, then $(M N) \in \Lam$.
  \item (Abstraction) If $u \in \Var$ and $M \in \Lam$, then
        $(\lambda u \,.\, M) \in \Lam$.
\end{enumerate}
The definition is \emph{inductive}: (1), (2) and (3) are the only ways to
construct elements of $\Lam$ (p.~4).
\end{definition}

\motivation{\pending{From \S1.2: functions as abstraction and application.}}

\begin{quantified}
\pending{Test 2 (least set): the intersection form. Name the domain the
outer quantifier ranges over. Then its negation: when is a string
\emph{not} in $\Lam$? Instantiate Definition~\ref{def:inductive-general}
(see correction~11 for how abstraction fits).}
\end{quantified}

\begin{convention}[Abstract syntax]\label{conv:abstract-grammar}
\uses{def:lambda-terms}
\source{N\&G \S1.3, p.~4 (\textbf{verify} number and page).}
$\Lam = \Var \mid (\Lam\Lam) \mid (\lambda \Var.\Lam)$.
\end{convention}

\begin{convention}[Syntactic identity]\label{conv:syntactic-identity}
\uses{def:lambda-terms}
\source{N\&G \S1.3 (\textbf{verify} number and page).}
Syntactic identity of $\lambda$-terms is written $\equiv$. For example,
$(xy) \equiv (xy)$ and $(xy) \not\equiv (yx)$.
\end{convention}

\begin{remark}
$\equiv$ is a metalanguage relation: identity of strings, Lean's
\texttt{Eq} on the inductive type of terms (CONVENTIONS \S7). Compare
\texttt{MetaEq} in \texttt{Aim/MetaLanguage/Eq.lean}.
\end{remark}

% ---------------------------------------------------------
\section{Where finiteness comes from}\label{sec:finiteness}

\begin{remark}[Inductive definitions: the pattern]\label{rem:inductive-pattern}
\uses{def:lambda-terms}
\source{Explanatory; general treatment in Enderton \S1.4 (Induction and
Recursion), pp.~34--44. Construction sequences first appear in Enderton
\S1.1.}

\paragraph{The pattern.} An inductive definition has three parts:
\begin{enumerate}[label=(\roman*)]
  \item a \emph{base}: a set of starting objects;
  \item \emph{rules}: finitely many operations, each taking finitely many
        objects and producing a new one;
  \item a \emph{closure clause}: nothing else is in the set.
\end{enumerate}
For $\Lam$: the base is $\Var$; the rules are application (two inputs) and
abstraction (a variable and one term); the closure clause is the sentence
after Definition~1.3.2. The same pattern defines $\mathbb{N}$ (base $0$, rule
successor), Enderton's $\WFF$ (base sentence symbols, rules the connectives),
derivations, typing derivations, and every Lean \texttt{inductive} type.

\paragraph{The rules recurse.} Rules (2) and (3) are recursive because they
refer to $\Lam$ itself: they take terms in and produce terms. There is no
separate recursion clause; the self-reference is the recursion.

\paragraph{The process never stops.} The rules apply again and again:
$(MN)$, $((MN)(MN))$, \dots{} are all terms. Arranged in stages,
\[
  \Lam_0 = \Var, \qquad
  \Lam_{n+1} = \Lam_n \cup \{(MN) : M,N \in \Lam_n\}
               \cup \{(\lambda u.M) : u \in \Var,\ M \in \Lam_n\},
  \qquad
  \Lam = \bigcup_{n} \Lam_n .
\]
\begin{center}
\fbox{There is no fixed bound on term size, but each individual term is
finite.}
\end{center}
Two separate facts, as in $\mathbb{N}$:
\begin{itemize}
  \item \emph{No largest element.} A term at stage $n$ has a larger one at
        stage $n+1$.
  \item \emph{No infinite element.} An endless expression such as
        $\lambda x.\lambda x.\lambda x.\cdots$ is not a member. This is not a
        limit on what can be chosen; the object is not in the set.
\end{itemize}
$\mathbb{N}$ has no largest number, and $\infty$ is not a number.

\paragraph{Where finiteness is enforced.} The closure clause is informal
(``the only ways to \emph{construct}''); the finiteness hides in the word
\emph{construct}. It is made precise in one of two ways, which define the
same set (Enderton \S1.4):
\begin{center}
\begin{tabular}{lll}
\toprule
Reading & Definition & Finiteness \\
\midrule
Bottom-up & last entry of a finite construction sequence & in the definition \\
Top-down & intersection of all sets closed under (1)--(3) & a theorem \\
\bottomrule
\end{tabular}
\end{center}
In the top-down reading three ingredients enforce finiteness: leastness, a
base of finite objects, and \emph{finitary} rules. The set of finite strings
is closed under the rules, and the least closed set lies inside every closed
set. A rule with infinitely many inputs would let infinite terms in.

\paragraph{Dual.} Taking the \emph{greatest} set consistent with the rules
instead (a coinductive definition) would admit infinite terms. ``Least'' is
exactly the choice that excludes them. Lean's \texttt{inductive} always means
the least.

\paragraph{Inductive, recursive, induction.}
\begin{center}
\begin{tabular}{lll}
\toprule
Notion & Does & Example \\
\midrule
Inductive definition & defines a set & $\Lam$ (Def.~1.3.2) \\
Recursive definition & defines a function on it, by cases on the rules
  & $\mathrm{Sub}$ (Def.~1.3.5) \\
Proof by induction & proves a property of every element, by cases
  & Lemma~1.3.6 \\
\bottomrule
\end{tabular}
\end{center}
Both of the last two rest on leastness: every element came from some rule.
Recursion needs one more hypothesis, \emph{free generation} (unique
readability): each term has exactly one last construction step. In Lean, one
object, the recursor, provides both recursion and induction.
\end{remark}

\begin{definition}[Inductive definition]\label{def:inductive-general}
\source{Enderton \S1.4, p.~34ff.\ (correction~9: Enderton's framework,
not N\&G's).}
Fix
\begin{enumerate}[label=(\roman*)]
  \item a set $U$, the \emph{universe};
  \item a set $B \subseteq U$, the \emph{base};
  \item a set $F$ of operations, each $f \in F$ a function
        $f\colon U^n \to U$ for some $n \ge 1$, its \emph{arity}.
\end{enumerate}
A set $S \subseteq U$ is \emph{closed} (under $B$ and $F$) if $B \subseteq S$
and for every $f \in F$ of arity $n$ and all $x_1,\dots,x_n \in S$,
$f(x_1,\dots,x_n) \in S$. The set \emph{inductively defined} by $B$ and $F$
is
\[
  C = \bigcap \{\, S \subseteq U : S \text{ is closed} \,\}.
\]
\end{definition}

\begin{quantified}
\[
x \in C \iff \forall S \subseteq U\,\bigl(S \text{ closed} \to x \in S\bigr)
\]
\end{quantified}

\begin{remark}
\pending{Why the family is nonempty (correction~10): which set is always
closed?}
\end{remark}

\begin{definition}[Construction sequence for $\Lam$]\label{def:construction-sequence-lambda}
\uses{ass:variables}
\source{Not in N\&G; Enderton's pattern (\S1.1, \S1.4) applied to $\Lam$.}
Let $A = \Var \cup \{(,\,),\,\lambda,\,.\}$, with $\frown$ for
concatenation of strings over $A$. A \emph{construction sequence} is a finite
sequence $\seq{t_0,\dots,t_n}$ of strings over $A$ such that for each
$i \le n$ at least one holds:
\begin{enumerate}[label=(\arabic*)]
  \item $t_i \in \Var$;
  \item $t_i = (\frown t_j \frown t_k \frown)$ for some $j,k < i$;
  \item $t_i = (\frown \lambda \frown u \frown . \frown t_j \frown)$ for some
        $u \in \Var$ and $j < i$.
\end{enumerate}
A string $t$ is \emph{constructible} if it is the last entry of some
construction sequence; $\Lam_*$ is the set of constructible strings.
\end{definition}

\begin{quantified}
\[
\begin{aligned}
t \in \Lam_* \iff \exists n\, \exists \seq{t_0,\dots,t_n}
\Bigl[\, & t_n = t \;\land\; \forall i \le n \bigl( t_i \in \Var \\
& \lor\; \exists j,k < i.\ t_i = (t_j t_k) \\
& \lor\; \exists u \in \Var\, \exists j < i.\ t_i = (\lambda u.\, t_j)
\bigr) \Bigr]
\end{aligned}
\]
Finiteness is explicit twice: in $\exists n$, and in $j,k < i$.
\end{quantified}

\begin{theorem}[Every $\lambda$-term is finite]\label{thm:lambda-terms-finite}
\uses{def:lambda-terms, rem:inductive-pattern}
\source{Explanatory; pattern from Enderton \S1.4.}
Every element of $\Lam$ is a finite string over $A$.
\end{theorem}

\begin{remark}[Proof structure]
Let $F$ be the set of finite strings over $A$. Check $F$ satisfies
(1)--(3); then leastness gives $\Lam \subseteq F$.
\end{remark}

\begin{proof}
\pending{Your proof.}
\end{proof}

\begin{remark}[Gaps]
\begin{itemize}
  \item N\&G assume, without proof, that the top-down and bottom-up readings
        agree. Proof: Enderton \S1.4. \pending{Your proof, in the Enderton
        notes.}
  \item N\&G define $\mathrm{Sub}$ by recursion on $\Lam$ without stating
        unique readability. \textbf{Verify} whether it is stated anywhere in
        Ch.~1.
\end{itemize}
\end{remark}

% ---------------------------------------------------------
\section{Subterms}

\begin{definition}[Multiset of subterms]\label{def:sub}
\uses{def:lambda-terms}
\source{N\&G Definition~1.3.5 (\textbf{verify} page).}
The multiset $\mathrm{Sub}(M)$ of subterms of $M$ is defined by recursion
on $M$:
\begin{enumerate}[label=(\arabic*)]
  \item (Basis) $\mathrm{Sub}(x) = \{x\}$ for each $x \in \Var$;
  \item (Application) $\mathrm{Sub}((MN)) = \mathrm{Sub}(M) \cup
        \mathrm{Sub}(N) \cup \{(MN)\}$;
  \item (Abstraction) $\mathrm{Sub}((\lambda x.M)) = \mathrm{Sub}(M) \cup
        \{(\lambda x.M)\}$.
\end{enumerate}
Here $\cup$ is \emph{multiset} union: a subterm may occur more than once
(correction~12). We call $L$ a \emph{subterm} of $M$ if
$L \in \mathrm{Sub}(M)$, i.e.\ $L$ occurs in $\mathrm{Sub}(M)$ at least once.
\end{definition}

\begin{example}
\pending{Compute $\mathrm{Sub}((xx))$ and say which element occurs twice.
Then a near miss: a string that occurs inside $M$ as a substring but is not
a subterm of $M$.}
\end{example}

\begin{remark}
$\mathrm{Sub}$ is the \emph{recursive definition} row of the table in
Remark~``Inductive definitions: the pattern''. That it is well defined
needs unique readability (see Gaps).
\end{remark}

\begin{lemma}[Properties of $\mathrm{Sub}$]\label{lem:sub-properties}
\uses{def:sub}
\source{N\&G Lemma~1.3.6.}
\begin{enumerate}[label=(\arabic*)]
  \item (Reflexivity) For all $\lambda$-terms $M$, $M \in \mathrm{Sub}(M)$.
  \item (Transitivity) If $L \in \mathrm{Sub}(M)$ and
        $M \in \mathrm{Sub}(N)$, then $L \in \mathrm{Sub}(N)$.
\end{enumerate}
\end{lemma}

\begin{quantified}
(correction~13)
\begin{align*}
&\forall M \in \Lam\ \bigl(M \in \mathrm{Sub}(M)\bigr)\\
&\forall L, M, N \in \Lam\ \bigl(L \in \mathrm{Sub}(M) \land
  M \in \mathrm{Sub}(N) \to L \in \mathrm{Sub}(N)\bigr)
\end{align*}
\end{quantified}

\begin{proof}
\pending{Your proofs. For each part: which variable do you induct on, and
what are the three cases?}
\end{proof}

\begin{definition}[Proper subterm]\label{def:proper-subterm}
\uses{def:sub, conv:syntactic-identity}
\source{N\&G \S1.3, Definition~1.3.7 (\textbf{verify} number).}
$L$ is a \emph{proper subterm} of $M$ if $L$ is a subterm of $M$ and
$L \not\equiv M$ (correction~14).
\end{definition}

\begin{quantified}
\pending{Predicate reading (test 6), and the remark that the relation
``is a proper subterm of'' is irreflexive (correction~15). Is it also
transitive?}
\end{quantified}

\begin{remark}[Lean]
\pending{After Enderton \S1.4: the general node (universe, base, operations;
$C^* = C_*$; induction) proposed for \texttt{Aim/MetaLanguage}, with $\Lam$,
$\WFF$ and $\mathbb{N}$ as instances; then the bridge showing a Lean
\texttt{inductive} gives the same set.}
\end{remark}

\end{document}
